\section{About the Committee}

\subsection{Mandate}

The Charter of the Security Council mandates it as the principal organ of the United Nations (UN) responsible for the maintenance of international peace and security.\footnote{(United Nations, 1945, Art. 24)} Article 24 of the UN Charter dictates that all UN Member States “agree to accept and carry out the decisions of the Security Council in accordance with the present Charter”. This means the Security Council is the only UN body which has the power to make legally binding decisions for all member states.\footnote{(United Nations, 1945, Art. 24-25)} Its responsibilities range from determining threats to international peace and security and recommending methods of peaceful resolution, to sanctioning and authorising the use of military force.\footnote{(United Nations, 1945, VI-VII)}

\subsubsection{Topics}

The UNSC deals with a broad range of issues, all connected to the core mandate of maintaining international peace and security. This includes:

\begin{itemize}
  \item \textbf{Threats to peace, breaches of the peace, and acts of aggression} -- Article 39 of the Charter allows the UNSC to address interstate wars, civil conflicts, and other politically destabilising situations.\footnote{United Nations. (1945). Charter of the United Nations, Chapter VII, Article 39. \url{https://www.un.org/en/about-us/un-charter/chapter-7}}
  \item \textbf{Interstate and intrastate conflicts} -- Development of the UN peacekeeping force, which deals with situations such as non-state armed groups involved in conflict and violent extremism.\footnote{United Nations. (n.d.). Peace and security. \url{https://www.un.org/en/global-issues/peace-and-security}}
  \item \textbf{Sanctions and economic measures} -- Includes various methods to pressure parties to comply with UNSC demands, without resorting to military force.\footnote{United Nations. (1945). Charter of the United Nations, Chapter VII, Article 41. \url{https://www.un.org/en/about-us/un-charter/chapter-7}} Methods include trade embargoes, arms embargoes, travel bans, freezing assets, and severance of diplomatic relations.\footnote{United Nations. (n.d.). Peace and security. \url{https://www.un.org/en/global-issues/peace-and-security}}
  \item \textbf{Peacekeeping and stabilisation} -- Operations which aim to facilitate political processes (such as elections or constitutional processes), promoting human rights, restoring the rule of law, and extending legitimate state authority.\footnote{United Nations. (n.d.). Maintain international peace and security. \url{https://www.un.org/en/our-work/maintain-international-peace-and-security}}
  \item \textbf{Terrorism} -- Coordinating global responses to fight against terrorism. The UNSC has also created subsidiary bodies dedicated to this cause (eg. counter-terrorism committees).\footnote{United Nations. (n.d.). Maintain international peace and security. \url{https://www.un.org/en/our-work/maintain-international-peace-and-security}}
  \item \textbf{Protection of civilians and prevention of atrocities} -- Responding to early warnings of genocide, war crimes, ethnic cleasing, and crimes against humanity. UNSC’s mandate is to protect civilian operations and is authorised to coordinate collective action if the national authorities are not able / do not.\footnote{United Nations. (n.d.). Security Council: Responsibility to protect. \url{https://www.un.org/en/genocide-prevention/responsibility-protect/security-council}}
  \item \textbf{Disarmament and arms regulations} -- Article 26 of the UN Charter allows for planning and creating regulation of armaments.\footnote{United Nations. (1945). Charter of the United Nations, Chapter V, Article 26. \url{https://www.un.org/en/about-us/un-charter/chapter-5}}
  \item \textbf{Dispute settlement and mediation} -- Engaging in preventative diplomacy, through mediation, investigations, appointing representatives, and requesting help from valuable stakeholders, before using enforcement measures.\footnote{United Nations. (n.d.). Peace and security. \url{https://www.un.org/en/global-issues/peace-and-security}}
\end{itemize}

\subsubsection{Composition}

The UNSC is composed of 15 members. There are two separate categories: five permanent members and ten non-permanent rotating members. Each member has one vote. The permanent members are China, France, the Russian Federation, the United Kingdom, and the United States. The non-permanent rotating members are elected by the UN General Assembly for two-year terms.\footnote{United Nations. (1945). Charter of the United Nations, Chapter V, Article 23. \url{https://www.un.org/en/about-us/un-charter/chapter-5}} Article 27 of the UN Charter requires 9 of 15 members to vote for procedural matters. Substantive matters require nine affirmative votes, but all the permanent members must vote in the affirmative, as they possess veto power.\footnote{United Nations. (n.d.). Main bodies. \url{https://www.un.org/en/about-us/main-bodies}} If any of the permanent members vote against, the substantive matter has been unilaterally stopped.\footnote{United Nations. (1945). Charter of the United Nations, Chapter V, Article 27. \url{https://www.un.org/en/about-us/un-charter/chapter-5}} Much academic scholarship has been conducted on this history of the veto power and how it has affected decision making within the Security Council.

\subsubsection{Operating Procedure}

The Security Council functions continuously, as intended in Article 28(1) of the UN Charter.\footnote{United Nations. (1945). Charter of the United Nations, Chapter V, Article 28. \url{https://www.un.org/en/about-us/un-charter/chapter-5}} Each member of the Council must be represented at all times, most often the UN headquarters in New York City. The Council also rotates the position of UNSC President, each month. The President is determined by an English alphabetical order of their country names (for both permanent and non-permanent members).\footnote{United Nations Security Council. (n.d.). Provisional rules of procedure (S/96/Rev.7), Rule 18. \url{https://main.un.org/securitycouncil/en/content/rop/chapter-4}} The President must call at least one meeting, but any Council member can request a meeting. The President must call a meeting when a dispute has been brought to the Council’s attention, for example, a situation regarding disturbance of international peace or security. The President is also obliged to call a meeting if the UN Secretary-General brings a matter to the Council’s Attention. An additional rule, allows for more formal semi-annual meetings periodically.\footnote{United Nations Security Council. (n.d.). Provisional rules of procedure (S/96/Rev.7), Rules 1–3. \url{https://main.un.org/securitycouncil/en/content/rop/chapter-1}} The Security Council meets much more than the minimum threshold. In 2025, the Security Council had 255 formal meetings and 115 informal consultations.\footnote{United Nations Security Council. (2025). Highlights 2025. \url{https://main.un.org/securitycouncil/en/content/highlights-2025}}

The UN Secretary-General draws up a provisional agenda for each Security Council meeting, and the President must approve it. Then the provisional agenda must be communicated to Council members before the meeting. Any unfinished agenda item from the last meeting, automatically carries over to the next meeting, unless the Council decides otherwise.\footnote{United Nations Security Council. (n.d.). Provisional rules of procedure (S/96/Rev.7), Rules 6–10. \url{https://main.un.org/securitycouncil/en/content/rop/chapter-2}}

The Security Council has many types of meetings available to it. Including formal meetings, which are the default type, and are normally public, either through transcript, meeting notes, or livestreaming.\footnote{United Nations Security Council. (n.d.). Repertoire of the Practice of the Security Council, Part II: Meetings and Records. \url{https://main.un.org/securitycouncil/sites/default/files/2025/26th\%20Supplement_part\%20II.pdf}} Private meetings are closed to the public.\footnote{United Nations Security Council. (n.d.). Repertoire of the Practice of the Security Council, Part II: Meetings and Records. \url{https://main.un.org/securitycouncil/sites/default/files/2025/26th\%20Supplement_part\%20II.pdf}} Informal consultations, also called consultations of the whole, bring all Council members into a private room with an agreed agenda, but without outside observers.\footnote{United Nations Security Council. (2025). Highlights 2025. \url{https://main.un.org/securitycouncil/en/content/highlights-2025}} \footnote{United Nations Security Council. (n.d.). Procedural issues. \url{https://main.un.org/securitycouncil/en/content/repertoire/procedural-issues}} Informal meetings are much more flexible and have been developed to depend more on the situation at the time.

Non-Council members may participate in the meeting if the Council members decide that their interests are “specially affected” by the agenda of the meeting.\footnote{United Nations. (1945). Charter of the United Nations, Chapter V, Article 31. \url{https://www.un.org/en/about-us/un-charter/chapter-5}} They do not have voting power on procedural or substantive matters. Any State, regardless of UN membership, which is a party to a dispute being reviewed by the Council must be invited to participate. They also do not have voting power.\footnote{United Nations. (1945). Charter of the United Nations, Chapter V, Article 32. \url{https://www.un.org/en/about-us/un-charter/chapter-5}}

When the Security Council receives a complaint about a threat to peace, it first recommends that the parties seek an agreement via peaceful means.\footnote{United Nations. (n.d.). Maintain international peace and security. \url{https://www.un.org/en/our-work/maintain-international-peace-and-security}} However, the Council has the power to investigate the issue, mediate between the parties, or use other diplomatic means to resolve the issue.\footnote{United Nations. (n.d.). Peace and security. \url{https://www.un.org/en/global-issues/peace-and-security}}

\subsubsection{Enforcement Powers}

The Security Council’s most influential power comes from enforcement. Chapter VII of the UN Charter governs the “Action with Respect to Threats to the Peace, Breaches of the Peace, and Acts of Aggression.”\footnote{United Nations. (1945). Charter of the United Nations, Chapter VII. \url{https://www.un.org/en/about-us/un-charter/chapter-7}} These include:

\begin{itemize}
  \item \textbf{Article 39} -- the Council may determine when a situation is a threat to the peace, a breach of the peace, or an act of aggression; and to decide on measures to maintain or restore international peace and security\footnote{United Nations. (1945). Charter of the United Nations, Chapter VII, Article 39. \url{https://www.un.org/en/about-us/un-charter/chapter-7}}
  \item \textbf{Article 40} -- the Council can call on parties to comply with measures to prevent situations from worsening\footnote{United Nations. (1945). Charter of the United Nations, Chapter VII, Article 40. \url{https://www.un.org/en/about-us/un-charter/chapter-7}}
  \item \textbf{Article 41} -- the Council can impose measures (without the use of armed force)\footnote{United Nations. (1945). Charter of the United Nations, Chapter VII, Article 41. \url{https://www.un.org/en/about-us/un-charter/chapter-7}} \footnote{United Nations. (n.d.). Uphold international law. \url{https://www.un.org/en/our-work/uphold-international-law}}
  \item \textbf{Article 42} -- the Council may (as a measure of last resort) take action by air, sea, or land forces if it is necessary to maintain or restore international peace and security\footnote{United Nations. (1945). Charter of the United Nations, Chapter VII, Article 42. \url{https://www.un.org/en/about-us/un-charter/chapter-7}}
\end{itemize}

\subsubsection{Self-Defense}

Article 51 of the UN Charter allows UN Member States the right of individual or collective self-defense.\footnote{United Nations. (1945). Charter of the United Nations, Chapter VII, Article 51. \url{https://www.un.org/en/about-us/un-charter/chapter-7}} The requirement is that an armed attack has been taken against them, and they may only act after the Security Council has taken enforcement measures. All acts of self-defense must be reported to the Council, and must not affect the Council’s authority or responsibility to act.\footnote{United Nations. (n.d.). Uphold international law. \url{https://www.un.org/en/our-work/uphold-international-law}}

\subsubsection{Peacekeeping Operations}

UN Peacekeeping operations follow mandates from the Security Council.\footnote{United Nations. (n.d.). Maintain international peace and security. \url{https://www.un.org/en/our-work/maintain-international-peace-and-security}} Peacekeeping operations are formed of troops and police, which are contributed by the Member States and are managed by the Department of Peace Operations.\footnote{United Nations. (n.d.). Maintain international peace and security. \url{https://www.un.org/en/our-work/maintain-international-peace-and-security}} More recently, peacekeeping has become multidimensional and goes beyond the original mandate of monitoring compliance with ceasefires. The peacekeeping forces may be called to help facilitate political processes, protect civilians, support disarmament, demobilisation, and reintegrating former combatants.\footnote{United Nations. (n.d.). Maintain international peace and security. \url{https://www.un.org/en/our-work/maintain-international-peace-and-security}}

\subsubsection{Subsidiary Organs}

Article 29 of the UN Charter allows the Security Council to establish subsidiary organs if it deems it necessary to perform its functions.\footnote{United Nations. (1945). Charter of the United Nations, Chapter V, Article 29. \url{https://www.un.org/en/about-us/un-charter/chapter-5}} This has included the creation of:

\begin{itemize}
  \item \textbf{Counter-Terorrism and Non-Proliferation committees} -- monitor implementation of counter-terorrism measures and measures to prevent non-states actors from acquiring nuclear, chemical, or biological weapons\footnote{United Nations Security Council. (n.d.). Counter-Terrorism Committee: Our mandate. \url{https://www.un.org/securitycouncil/ctc/content/our-mandate}}\footnote{43}
  \item \textbf{Sanction Committees} -- operate regionally to monitor the effect and compliance with economic sanctions\footnote{United Nations Security Council. (n.d.). Sanctions. \url{https://main.un.org/securitycouncil/en/sanctions/information}} \footnote{United Nations Security Council. (n.d.). Security Council subsidiary bodies: An overview. \url{https://main.un.org/securitycouncil/en/content/security-council-subsidiary-bodies-overview}}\footnote{United Nations Security Council. (n.d.). Counter-Terrorism Committee: FAQs. \url{https://www.un.org/securitycouncil/ctc/content/faqs}}\footnote{United Nations Security Council. (n.d.). 1540 Committee: Frequently asked questions. \url{https://www.un.org/en/sc/1540/faq.shtml}}
  \item \textbf{Military Staff Committee} -- advises and assists the Security Council on military requirements for its functions, regulation of armaments, and the employment and command of armed forces\footnote{United Nations Security Council. (n.d.). Security Council subsidiary bodies: An overview. \url{https://main.un.org/securitycouncil/en/content/security-council-subsidiary-bodies-overview}}
  \item \textbf{Standing Committees and Ad Hoc Bodies} -- addressing procedural matters (eg. applications for UN Membership)\footnote{United Nations. (n.d.). Repertoire of the Practice of the Security Council, Supplement 2000–2003, Chapter V: Subsidiary Organs. \url{https://www.un.org/french/docs/cs/repertoire/00-03/00-03_5.pdf}}
\end{itemize}

\subsubsection{Peacekeeping Operations and Special Political Missions}

\begin{itemize}
  \item \textbf{International Courts and Tribunals} -- to complete the mandate of International Criminal Tribunals, and to adjudicate criminal charges against larger regimes\footnote{United Nations Security Council. (n.d.). Security Council subsidiary bodies: An overview. \url{https://main.un.org/securitycouncil/en/content/security-council-subsidiary-bodies-overview}}
  \item \textbf{Peacebuilding Commission} -- intergovernmental advisory body to support peace efforts in countries coming from conflict; advises both the Security Council and the UN General Assembly\footnote{United Nations. (n.d.). Peacebuilding. \url{https://www.un.org/peacebuilding/}}
\end{itemize}
